\section{Scope of this slice}
\label{sec:tech:scope}

The full system has two halves. Private \emph{training} learns a
factorisation of a rating matrix that no server may see, and private
\emph{delivery} fetches the recommended item without revealing which item it
was. The requirements document fixes the build order as S1, then S2, then
S3: serving and delivery first, then training, then the composition. S1 is
first because it is lower risk, it demonstrates something end to end early,
and the distributed point function it needs is a prerequisite for the
non-linear gates that private training will need later.

This report covers S1 only. The model is trained \emph{in the clear} by the
Python oracle of Section~\ref{sec:tech:oracle}, and everything downstream of
that model is private. Private training is the subject of the next phase and
no claim is made about it here.

\subsection{What is deliberately absent}
\label{sec:tech:scope:absent}

Three components specified in the design draft are not built in this slice,
and each was cut for a stated reason rather than for lack of time.

A distributed comparison function is not implemented. Its only consumers are
the truncation and normalisation protocols, both of which belong to private
training. The design draft asserts that a distributed point function and a
distributed comparison function are the same primitive. They are not, and the
reference implementation of \nudgeplain{} keeps them in separate modules,
which corroborates the distinction.

Oblivious top-$\topk$ selection is not implemented. The user reconstructs the
score vector and selects the top $\topk$ locally, so no server learns the
selection regardless of whether an oblivious selector exists. The apparatus
would protect something that is already protected.

Network emulation is not used. The development machine has neither Docker nor
\texttt{tc netem}, so every measurement in this report is taken on the
\texttt{local} profile and is labelled as such. Wide-area numbers are
deferred rather than estimated.

Batching is not implemented. \textsc{Architecture} \S8 gives the network layer
``framing, batching, \texttt{flush()}'', and the first two words are built. But
this slice exchanges a handful of frames per query, so there is nothing to
batch; \texttt{flush()} exists and is documented as the no-op it honestly is
under blocking sockets with no user-space buffer. An unused optimisation would
be harder to justify than its absence.

Wide-area behaviour is not measured, for the same reason network emulation is
not used. Every number here is on the \texttt{local} profile and says so.

\medskip

\noindent
An earlier version of this section said networking itself was not implemented
and named it as the largest open gap in the slice. That was true when written
and is no longer: the three servers are now three separate operating-system
processes talking over TCP, and Section~\ref{sec:tech:net} reports the
transcript and distinguisher experiment that the phase's exit criterion asks
for.
